\section{Comparative Analysis}
\label{sec_comparative_analysis}

\reftab{tab:vapt-comparison} summarizes a representative subset of the surveyed systems along four dimensions that recurred throughout the literature review: whether the system uses a single agent or a coordinated team, whether it models dependencies between candidate weaknesses explicitly, whether it retains any form of memory across sessions, and which benchmark(s) it was evaluated against. The full survey covers more systems than shown here; the table is restricted to the systems most directly comparable to one another.

\begin{table}[h!]
\centering
\begin{tabular}{p{2.4cm}p{2.1cm}p{2.6cm}p{2.1cm}p{3cm}}
\hline
\textbf{System} & \textbf{Architecture} & \textbf{Dependency Modeling} & \textbf{Cross-Session Memory} & \textbf{Benchmark(s)} \\
\hline
Fang et al.~\cite{fang2024oneday} & Single agent (ReAct) & None (given CVE description) & None & 15 one-day CVEs \\
PentestGPT~\cite{deng2024pentestgpt} & Single agent, split reasoning/parsing & Implicit (LLM reasoning only) & None & HTB/VulnHub machines \\
HPTSA~\cite{zhu2024zeroday} & Hierarchical planner + task agents & None (flat dispatch) & None & 15 zero-day CVEs \\
MAPTA~\cite{david2024mapta} & Multi-agent team & None (flat dispatch) & None & CTF challenges \\
VulnBot~\cite{kong2025vulnbot} & Multi-agent, role-specialized & None (flat dispatch) & None & HTB-style scenarios \\
CHECKMATE~\cite{wang2025classicalplanning} & Agent + classical planner & Explicit, pre-enumerated & None & Curated scenarios \\
PentestEval SMP~\cite{yang2025pentesteval} & Modular pipeline & Explicit, ground-truth-injected & None & 12 real-world scenarios \\
Incalmo~\cite{singer2025incalmo} & Multi-host orchestration & Partial (host/credential graph) & None & MHBench (40 environments) \\
PrediQL~\cite{liu2025prediql} & LLM-guided fuzzer & Schema-derived & None & 6 GraphQL APIs \\
\hline
\end{tabular}
\caption{Comparison of representative autonomous VAPT systems surveyed}
\label{tab:vapt-comparison}
\end{table}

A few patterns are already visible from this comparison, even at this early stage of the review. First, systems that scale well to a wide, unfamiliar attack surface (HPTSA, MAPTA, VulnBot) tend to dispatch tasks flatly, without an explicit model of which weaknesses depend on which others. Second, systems that do model dependencies explicitly (CHECKMATE, PentestEval's SMP configuration) tend to be evaluated on curated scenarios where the relevant weaknesses are already known in advance, rather than on open-ended discovery. Third, cross-session memory -- the ability for an agent to retain and reuse what it learned in a previous session against a similar target -- is largely absent from the systems surveyed so far, with the exception of ideas borrowed from non-security agent work such as VOYAGER~\cite{wang2023voyager} and Reflexion~\cite{shinn2023reflexion}. Finally, every system in the table is evaluated against a single benchmark family (web, GraphQL, or multi-host), and none of the surveyed papers report results across more than one attack-surface category using the same architecture.

These observations are preliminary and are discussed further, together with their implications, in the following section.
